\chapter{Fundamentos de finanzas corporativas}\label{ch:corporate-finance-core}

% Alcance: valor del dinero en el tiempo, NPV, la tasa de descuento como hurdle rate, la perpetuidad creciente.
% Referencia canónica del estilo de la edición: prosa fluida bajo títulos de sección, con
% números trabajados en `example` y advertencias precisas en `admonition`.
% El gemelo de marketing de la perpetuidad creciente se encuentra en \cref{ch:customer-economics}.

\section{La perpetuidad creciente}\label{sec:growing-perpetuity}
% complexity: 3

¿Cuánto vale hoy un flujo de caja que crece a una tasa constante \(g\) de forma indefinida? La
pregunta no es académica: este modelo único subyace al valor terminal de los flujos de caja
descontados y al valor de vida del cliente por igual, y un error en él se propaga a través de
toda valoración que dependa de ambos. La intuición es el fundamento de toda la disciplina
financiera. Un dólar el año próximo vale menos que un dólar hoy, porque el dólar de hoy puede
invertirse; el crecimiento compensa ese descuento solo en parte. Cuando la tasa de descuento
supera la tasa de crecimiento, una corriente infinita de flujos crecientes suma de todos modos
un número finito.

Sumando los flujos de caja descontados y reconociendo una serie geométrica de razón
\(\frac{1+g}{1+r}<1\), la suma infinita se colapsa en una sola razón:
\begin{equation}\label{eq:gordon}
  PV \;=\; \sum_{t=1}^{\infty}\frac{D\,(1+g)^{t-1}}{(1+r)^{t}}
      \;=\; \frac{D}{\,r-g\,}.
\end{equation}
El denominador \((r-g)\) es la fricción neta: la presión del descuento menos el impulso del
crecimiento. La serie converge \emph{solo} cuando \(r>g\). Si \(g\ge r\), el denominador
se vuelve cero o negativo y \cref{eq:gordon} devuelve un resultado sin sentido --- el error
de valoración más común, y el primero que un lector escéptico debería verificar en cualquier
modelo que reciba.

\begin{example}
Una cohorte de suscriptores genera \money{1250000} el año próximo con un margen de contribución
del \qty{42}{\percent}, que crece \qty{3}{\percent} por año, y la empresa descuenta al \qty{11}{\percent}. El
margen es lo que efectivamente acumula la empresa, por lo que es la \(D\) en \cref{eq:gordon};
el \(\qty{11}{\percent}\) es el costo de oportunidad del capital inmovilizado en atender la
cohorte:
\[
  PV \;=\; \frac{\money{1250000}\times 0.42}{0.11-0.03}\;\approx\;\money{6562500}.
\]
La cohorte vale aproximadamente \money{6562500} hoy. Nótese cuán sensible es esa cifra a la
brecha \(r-g\): ampliarla al \qty{10}{\percent} reduce el valor en una quinta parte, razón por la
cual el supuesto de crecimiento merece más escrutinio que cualquier otro insumo del modelo.
\end{example}

\begin{admonition}{Advertencia}
Reutilizar la misma \(g\) elevada para diez años de pronóstico explícito \emph{y} para el valor
terminal duplica el crecimiento. La \(g\) terminal debe reflejar únicamente el estado estacionario
--- generalmente igual o inferior al crecimiento del PIB a largo plazo. Una tasa terminal tomada
del hipercrecimiento reciente de una empresa es la forma más rápida de fabricar una valoración
que no resiste el contacto con la realidad.
\end{admonition}

Al relajar el crecimiento constante, el modelo se convierte en uno de dos etapas o en un
modelo H; al relajar el \emph{churn} constante en su gemelo de \emph{marketing}, el valor del
cliente requiere un modelo de supervivencia \emph{beta}-geométrica desplazada. La fórmula
limpia es el \emph{piso} del análisis, no el techo. El álgebra idéntica reaparece como valor
de vida del cliente en \cref{ch:customer-economics}: renombrando \(D\) como el margen anual por
cliente y \(g\) como su tasa de crecimiento, \cref{eq:gordon} valora a un cliente exactamente
igual que valora una cohorte. \textcite{brealey2020} desarrollan en profundidad el lado de las
finanzas corporativas.

\section{El valor del dinero en el tiempo}\label{sec:time-value-of-money}
% complexity: 2

¿Cuánto vale hoy un pago futuro en efectivo? Toda decisión de asignación de capital ---
invertir, adquirir o recomprar acciones --- depende de responder esto antes que cualquier
otra cosa. Un dólar en mano hoy puede reinvertirse a la tasa \(r\), acumulando \((1+r)^n\)
dólares en \(n\) años; recorriendo esa lógica en sentido inverso, un pago futuro vale menos
que su valor nominal exactamente por el factor que el interés compuesto habría producido.

Si una suma \(\mathrm{FV}\) llega en \(n\) años y el costo de oportunidad del capital es \(r\)
por período, su valor presente se obtiene descontando:
\begin{equation}\label{eq:pv}
  PV \;=\; \frac{\mathrm{FV}}{(1+r)^{n}}.
\end{equation}
El denominador \((1+r)^{n}\) es el factor de valor futuro; su recíproco es el
\emph{factor de descuento}. \Cref{eq:gordon} no es más que la aplicación repetida de \cref{eq:pv}
a lo largo de una corriente infinita y creciente. La fórmula supone que \(r\) es constante y
conocida, que el \emph{cash flow} llega en la fecha indicada y que el inversor puede
efectivamente reinvertir al mismo \(r\); para horizontes cortos o flujos de bajo riesgo el error
derivado de estos supuestos es pequeño, pero para flujos de caja ilíquidos y de largo plazo
puede dominar el resultado.

\begin{example}
¿Cuánto vale hoy un pago de \money{252} que llega en tres años, a la tasa de descuento del
\qty{11}{\percent} de la empresa? Esa tasa es el costo de oportunidad del capital propio derivado
en \cref{ch:risk-and-capital-structure}; aplicando \cref{eq:pv},
\[
  PV \;=\; \frac{\money{252}}{(1.11)^{3}}
      \;=\; \frac{\money{252}}{1.3676}
      \;\approx\; \money{184}.
\]
La brecha de \money{68} entre el valor nominal y el valor presente es el costo de esperar ---
no la inflación, sino el interés compuesto no percibido. La decisión que esto informa: una
promesa de pagar \money{252} en tres años vale la pena aceptarla solo si el uso alternativo de
\money{184} hoy rinde menos del \qty{11}{\percent}.
\end{example}

\begin{admonition}{Advertencia}
Descontar flujos nominales con una tasa real, o viceversa, es el error clásico. La consistencia
entre el flujo y la tasa es el único invariante de la fórmula de valor presente: los flujos
nominales requieren una \(r\) nominal; los flujos ajustados por inflación requieren una \(r\)
real. Mezclarlos distorsiona el valor aproximadamente en la tasa de inflación por año, que se
acumula silenciosamente a lo largo de un proyecto de varios años.
\end{admonition}

Para flujos de caja riesgosos, la \(r\) «correcta» es en sí misma una salida del modelo, no un
dato --- véase \cref{eq:capm} en \cref{ch:risk-and-capital-structure}. La fórmula de valor
presente es la unidad atómica de todo el trabajo de valoración en este libro; \textcite{brealey2020}
la derivan desde primeros principios y la extienden a bonos, anualidades y el caso de la
perpetuidad creciente de \cref{eq:gordon}.

\section{Valor presente neto}\label{sec:net-present-value}
% complexity: 3

¿Debe la empresa realizar una inversión? El valor presente neto convierte un calendario
completo de flujos de caja --- un desembolso hoy, retornos a lo largo del tiempo --- en un
número único cuyo signo resuelve la pregunta. Un NPV positivo significa que la inversión rinde
más que la \emph{hurdle rate}; uno negativo significa que destruye valor a esa tasa. La
pregunta de fondo es directa: ¿cuánto más rico hace esto al inversor, medido en dólares de hoy,
después de cargar completamente el costo del capital desplegado?

Sea \(\mathrm{FCF}_t\) el \emph{cash flow} libre en el período \(t\), \(I_0\) el desembolso
inicial en \(t=0\), y \(r\) la tasa de descuento:
\begin{equation}\label{eq:npv}
  \mathrm{NPV} \;=\; \sum_{t=1}^{T}\frac{\mathrm{FCF}_t}{(1+r)^{t}} \;-\; I_0.
\end{equation}
Cada término \(\mathrm{FCF}_t/(1+r)^t\) es \cref{eq:pv} aplicada al período \(t\). La tasa de
descuento \(r\) es el costo de capital promedio ponderado de la empresa, derivado en
\cref{ch:risk-and-capital-structure} como \cref{eq:wacc}: captura tanto el costo del capital
propio (\cref{eq:capm}) como el costo de la deuda con su escudo fiscal, ponderados por la
estructura de capital. La fórmula supone que los flujos de caja son conocidos, \(r\) es
constante y el proyecto no incorpora opcionalidad --- ninguna capacidad de expandirse, abandonarse
o diferirse --- por lo que puede valorar erróneamente inversiones cuyo principal valor reside en
la flexibilidad que crean.

\begin{example}
¿Adquirir un cliente a \money{600} crea valor a la \emph{hurdle rate} del \qty{11}{\percent} de la
empresa? Tratar el costo de adquisición de \money{600} como el desembolso \(I_0\). El margen de
contribución anual del cliente de \money{252} crece al \qty{3}{\percent}, de modo que el flujo de
entradas es una perpetuidad creciente, valorada por \cref{eq:gordon}:
\[
  \sum_{t=1}^{\infty}\frac{\mathrm{FCF}_t}{(1.11)^t}
  \;=\; \frac{\money{252}}{0.11-0.03}
  \;=\; \money{3150}.
\]
Aplicando \cref{eq:npv},
\[
  \mathrm{NPV} \;=\; \money{3150} \;-\; \money{600} \;=\; \money{2550}.
\]
Cada dólar de costo de adquisición devuelve \money{5.25} de valor presente --- un excedente muy
por encima de la \emph{hurdle rate}. La implicación es que la restricción vinculante aquí no es
el costo del capital sino la oferta de clientes: la empresa debería aumentar su presupuesto de
adquisición hasta que el costo marginal de un cliente se aproxime a los \money{3150} que ese
cliente vale.
\end{example}

\begin{admonition}{Advertencia}
Usar una tasa de descuento demasiado baja para fabricar un NPV positivo es el pecado cardinal
de la evaluación de inversiones. Un proyecto que parece atractivo al \qty{6}{\percent} pero es
negativo al correcto \qty{11}{\percent} no crea valor --- lo destruye lentamente, hasta que el
verdadero costo del capital se impone a través de presión de refinanciamiento o dilución del
capital propio. Anclar la \emph{hurdle rate} al costo real del capital es la disciplina que
\cref{eq:npv} existe para hacer cumplir.
\end{admonition}

El NPV es el análogo corporativo de la decisión de economía del cliente en
\cref{ch:customer-economics}: ambos preguntan si un desembolso presente --- \(I_0\) o el costo
de adquisición --- se justifica por la corriente descontada que genera. La equivalencia no es
una coincidencia; es la lógica unificadora de la columna vertebral de creación de valor.
\textcite{brealey2020} demuestran la consistencia del NPV con la maximización del valor para
el accionista, y \textcite{ross2022} analizan su comparación con la tasa interna de retorno y
las condiciones bajo las cuales ambos divergen.
